\documentclass[11pt,a4paper]{article}
\usepackage[margin=2.5cm]{geometry}
\usepackage{amsmath,amssymb,amsfonts,graphicx,float,booktabs,xcolor,xparse,enumitem,soul}
\usepackage[hidelinks]{hyperref}
\setlength{\parskip}{0.6em}\setlength{\parindent}{0pt}
\definecolor{newc}{HTML}{1f6f43}
\newcommand{\new}[1]{\textcolor{newc}{#1}}
\NewDocumentCommand{\myfig}{O{htbp} m m m O{\linewidth}}{\begin{figure}[#1]\centering\includegraphics[width=#5]{#2}\caption{#3}\label{fig:#4}\end{figure}}
\newcommand{\citep}[1]{\textup{[#1]}}
\newcommand{\nref}[1]{\eqref{#1}}\newcommand{\Nref}[1]{\eqref{#1}}
\renewcommand{\thesubsection}{V-\Alph{subsection}}
\setcounter{subsection}{2}
\setcounter{equation}{4}
\title{Proposed changes to ``Where the items are'', version 3}
\date{10 November 2026}
\begin{document}
\maketitle

\section*{How to read this}
Part~1 lists the changes in plain words. Part~2 shows the Problem
Formulation section as it would read. Part~3 shows what changes in the
Implementation chapter and Appendix~C. Text in \new{green} is proposed and not
in the thesis yet, and \st{struck} text would go. Everything in black is
already in the thesis file. Part~4 lists what the check of the code found
and the questions for you.

\section*{Part 1. The changes}
\subsection*{The idea}
The problem is stated the way Waubert de Puiseau et al.\ (2022) state theirs.
It has a decision, what may be seen, what may not be done, and a cost. How
the problem is solved is separate. Here it is solved by three rules, and a
rule does not break the storage location assignment problem. It is one way
to solve it, as are the rule benchmarks in their paper and the classic
slotting rule. The only thing a rule may not claim is to be the best
possible placement.

The real cost is the mean order service time, the main measure of the
thesis. Until now the section gave $D + \beta K$ as the objective, but that
is only the stand-in the rules rank shelves by. So $D$, $K$ and $\beta$ move
to the Implementation chapter, where the rules are.

\subsection*{Changes in the Problem Formulation (Part 2)}
\begin{itemize}[leftmargin=1.2em]
\item \textbf{How much arrives.} New sentences say that stock comes in
refills of $Q$ units, that a refill goes whole onto one shelf, and that a SKU
gets one refill for every $Q$ units ordered. This was missing, and it is
what keeps the arrivals the same under every rule.
\item \textbf{A new example and Figure~\ref{fig:storageupdate}}, with three
refills of 3 units of tea. Two are placed and one fits nowhere, which also
shows why condition \eqref{eq:upd:arrived} uses $\le$.
\item \textbf{``What the rule decides'' becomes ``The storage problem''}. One
sentence names the problem and its four parts (decision, what is seen, the
conditions, the cost) without repeating them. A second paragraph names the
three rules in one sentence each and says plainly that congestion-aware
storage is the thesis's own, because the DSLAP studies have no traffic.
\item \textbf{$D$, $K$, $\beta$ and the objective leave this section.} They
were put here earlier today and would now move (Part~3). $\alpha$ stays
dropped.
\end{itemize}
\subsection*{Changes in the Implementation chapter (Part 3)}
\begin{itemize}[leftmargin=1.2em]
\item \textbf{The stand-in cost} $D + \beta K$ moves here as a new paragraph,
with the label \texttt{eq:storageobjective}. $D$ now follows the code. Each
leg of the trip counts as often as it is made, the pick leg once per order
line and the put-away leg once per refill. The sentence that gave both legs
the same weight goes.
\item \textbf{Ties} go by the trip distance, then by index, as in the code.
\item \textbf{Appendix~C's arrivals} are rewritten to the refill rule the code
uses, with a new symbol $\nu_t(k)$ for the units ordered since the last
refill.
\item \textbf{Other places} name the crowding cost in words, not as an
equation of a later chapter.
\item \textbf{Whole refills.} The sentence ``as many on each shelf as fit''
becomes ``each refill whole on the first shelf with room for it'', as the
code does.
\end{itemize}

\clearpage
\section*{Part 2. The Problem Formulation section as it would read}
\subsection{Where the items are}
\label{sec:pf:storage}

Stock is kept on shelves, and an item can only be picked from a shelf that
holds it. So where the stock sits decides where the agents must move to.
The shelves never move. What changes over time is what each shelf holds.

The storage state is a counting function. It says how many units of each
SKU sit on each shelf, for example 12 units of tea on shelf $A$. A shelf
can hold several SKUs, and one SKU can sit on several shelves.

Formally, $x_t(k,v)$ is the number of units of SKU $k$ on shelf $v$ at
time $t$. Here $k$ is one of the SKUs in $\mathcal{K}$, and $v$ one of the
shelves in $V_{\mathrm{str}}$:
\begin{equation}
x_t:\mathcal{K}\times V_{\mathrm{str}}\to\mathbb{N}_0 .
\label{eq:storagestate}
\end{equation}

Each shelf has limited room, and units of different SKUs take different
amounts of it. The units on a shelf, each counted with the room it takes,
may not need more room than the shelf has. Formally, let $b_k$ be the room
one unit of SKU $k$ takes and $\mathrm{cap}(v)$ the room on shelf $v$. The
storage states that fit are
\begin{equation}
\mathcal{X}=
\left\{
x\in\mathbb{N}_0^{\mathcal{K}\times V_{\mathrm{str}}}:
\sum_{k\in \mathcal{K}}b_kx(k,v)\le\mathrm{cap}(v)
\quad\forall v\in V_{\mathrm{str}}
\right\}.
\label{eq:feasiblestorage}
\end{equation}

\paragraph{Example.}
Table~\ref{tab:storagematrix} shows a small storage state. Tea and mugs
take one unit of room each, and coffee takes two. Shelves $A$, $B$ and $C$
have room for 20, 15 and 10. Shelf $A$ is then exactly full, and $B$ and
$C$ have room left:
\[
\begin{pmatrix}
1&2&1
\end{pmatrix}
\begin{pmatrix}
12&0&5\\
4&3&0\\
0&6&2
\end{pmatrix}
=
\begin{pmatrix}
20&12&7
\end{pmatrix}
\le
\begin{pmatrix}
20&15&10
\end{pmatrix}.
\]
Room also ties SKUs together. Tea and coffee share the room on shelf $A$,
so more tea there leaves less room for coffee.

\begin{table}[H]
\centering
\caption{A small storage state.}
\label{tab:storagematrix}
\begin{tabular}{lccc}
\toprule
 & \(A\) & \(B\) & \(C\)\\
\midrule
tea & 12 & 0 & 5\\
coffee & 4 & 3 & 0\\
mugs & 0 & 6 & 2\\
\bottomrule
\end{tabular}
\end{table}

\paragraph{How stock changes.}
Stock leaves the shelves when agents pick it for orders, and new stock
comes in to refill them.

New stock arrives at a dock every $\Delta$ timesteps. Each of these moments
is a storage update. \new{Stock arrives in refills of $Q$ units, and a refill
always goes whole onto one shelf. How much arrives depends on the orders
alone. A SKU gets one refill for every $Q$ units ordered since its last
refill, so the arrivals are the same under every rule.} At a storage update, the storage rule $F$ chooses a
shelf for each arriving unit. An agent then carries the units from the dock
to that shelf. This trip is called a put-away.

While the units are on their way, they already take up room on the shelf,
but no order can pick them yet. So they are counted on their own. Write
$x^{\mathrm{in}}_t(k,v)$ for the units of SKU $k$ on their way to shelf
$v$, with ``in'' for inbound. When the agent delivers them, they count as
on the shelf.

\paragraph{What the rule knows.}
The rule only knows the last storage period. It reads three averages per
step over that period:
\begin{itemize}
\item \(\hat\rho_t(k)\), how often SKU \(k\) was ordered
\item \(\hat\mu_t(e)\), how often agents crossed edge \(e\)
\item \(\hat w_t(e)\), how many steps agents waited to get onto edge \(e\)
\end{itemize}
Demand storage uses only the first, the orders. Congestion-aware storage
uses all three. The last two come from how the agents moved, and show which
aisles were crowded. An edge counts as crowded when many agents cross it
and they often wait to get onto it.

\paragraph{The storage update.}
At a storage update the rule only decides where the new units go, and
three conditions must hold. No shelf may hold more than it has room for,
counting the units on their way. Units already on their way keep the shelf
they were given. And the rule cannot place more units than arrived, so
units that fit on no shelf are left unplaced.

Formally, the rule takes the stock just before the update, $x_t^{-}$ on the
shelves and $x^{\mathrm{in},-}_t$ on its way. It also takes $y_t(k)$, the
units of each SKU $k$ that arrived, and the three averages above. It
returns where the units go:
\begin{equation}
x_t^{\mathrm{in}}
=
F\left(
x_t^{-},\,
x_t^{\mathrm{in},-},\,
y_t,\,
\hat\rho_t,\,
\hat\mu_t,\,
\hat w_t
\right).
\label{eq:storageupdate}
\end{equation}
The three conditions, in the same order, are
\begin{align}
x_t^{-}+x_t^{\mathrm{in}}&\in \mathcal{X}
&&\text{no shelf over its room,}
\label{eq:upd:room}\\
x_t^{\mathrm{in}}&\ge x_t^{\mathrm{in},-}
&&\text{units on their way keep their shelf,}
\label{eq:upd:keep}\\
\sum_{v\in V_{\mathrm{str}}}
\left(
x_t^{\mathrm{in}}(k,v)-x_t^{\mathrm{in},-}(k,v)
\right)
&\le y_t(k)\quad\forall k\in \mathcal{K}
&&\text{no more placed than arrived.}
\label{eq:upd:arrived}
\end{align}

\paragraph{Example.}
\new{Take the stock of Table~\ref{tab:storagematrix} just before an update,
with nothing on its way, and say three refills of 3 units of tea arrive.
Shelf $A$ is full, and $B$ and $C$ each have room for one refill, so the
rule puts one refill on $B$ and one on $C$ (Figure~\ref{fig:storageupdate}).
Every shelf stays within its room, so \eqref{eq:upd:room} holds. Nothing was
on its way before, so \eqref{eq:upd:keep} holds. The third refill fits
nowhere and is left unplaced. The rule places 6 of the 9 units, so
\eqref{eq:upd:arrived} holds with $6 \le 9$.}

\myfig[H]{storage_update.png}{\new{One storage update on the stock of Table~\ref{tab:storagematrix}. Two refills of tea go whole to shelves with room, and the third fits nowhere.}}{storageupdate}[0.85\linewidth]

\paragraph{\new{The storage problem.}}
\new{Choosing a shelf when stock arrives is the dynamic storage location
assignment problem (Waubert de Puiseau et al., 2022). The rule makes the
storage update \eqref{eq:storageupdate}, sees only the past, must keep
conditions \eqref{eq:upd:room} to \eqref{eq:upd:arrived}, and is judged by
the mean order service time of Section V-J.}

\new{That cost depends on how the agents move after the update, so no rule
can work it out when it decides. Each rule therefore ranks the shelves by a
simple stand-in for it, given in Section VI-A. Fixed storage sends every
refill to its SKU's home shelf. Demand storage puts popular SKUs where their
trip is short, the classic slotting rule (Heskett, 1963; de Koster et al.,
2007). Congestion-aware storage also avoids shelves whose trip runs through
crowded aisles. The dynamic storage studies have no traffic, so this last
rule is this thesis's own, and it is the one that closes the loop.}

\paragraph{Example.}
Figure~\ref{fig:storageexample} shows tea that has become popular. Under
fixed storage, new tea goes back to its home shelf, three edges from the
workstation. Under demand storage, it goes to the nearest shelf with room,
one edge away. The old tea stays where it is and runs down as it is picked.

\myfig[t]{arrival_storage.png}{Arriving tea under fixed and demand storage. Old tea stays where it is, so only the new tea's shelf differs.}{storageexample}[0.8\linewidth]



\clearpage
\section*{Part 3. The Implementation chapter, the part that changes}
\subsection*{VI-A Storage rules (excerpt)}
All three rules place arriving stock the same way, and only their order
differs. They take one SKU at a time and put its arriving
\st{units on the shelves in the rule's order, as many on each shelf as fit}
\new{refills on the shelves in the rule's order, each refill whole on the first shelf with room for it}.
What fits nowhere is not placed. Stock already on a shelf stays where it is.

\textit{(The trip distance paragraph, fixed storage and the start of demand storage stay as they are.)}

\textit{In the congestion-aware paragraph:}
\st{Both legs carry the same $\hat\rho_t(k)$, because arrivals refill what
was ordered, so a SKU is put away about as often as it is picked.}

\paragraph{\new{The stand-in cost.}}
\new{The rules cannot see the service time they are judged by, so they rank
shelves by a stand-in. The classic slotting cost is each trip's length times
how often the trip is made (Heskett, 1963; de Koster et al., 2007). A shelf
has two trips. The pick leg runs from the shelf to a workstation of the
SKU's section and is made about $\hat\rho_t(k)$ times per step, once per
order line. The put-away leg runs from a dock to the shelf and is made about
$\hat\rho^{\mathrm{unit}}_t(k)/Q$ times per step, because one refill brings
$Q$ units.}

\new{Formally, let $d^{\mathrm{pick}}_t(k,v)$ be the mean distance from shelf
$v$ to the workstations of section $\hat h_t(k)$, $d^{\mathrm{put}}(v)$ the
mean distance from the docks to $v$, and $x^{\mathrm{new}}_t = x^{\mathrm{in}}_t - x^{\mathrm{in},-}_t$
the units placed at an update. The slotting cost is}
\begin{equation*}
\new{D = \sum_{k\in\mathcal{K}} \sum_{v\in V_{\mathrm{str}}}
\Bigl(\hat\rho_t(k)\, d^{\mathrm{pick}}_t(k,v) + \frac{\hat\rho^{\mathrm{unit}}_t(k)}{Q}\, d^{\mathrm{put}}(v)\Bigr)\, x^{\mathrm{new}}_t(k,v) .}
\end{equation*}
\new{The crowding cost is this thesis's own. It adds up the crowding
$\hat\mu_t(e)\,\hat w_t(e)$ on every edge $e$ of the trip through $v$, from
the docks to $v$ and on to the nearest workstation,}
\begin{equation*}
\new{K = \sum_{k\in\mathcal{K}} \sum_{v\in V_{\mathrm{str}}}
\Bigl(\sum_{e \in \mathrm{path}(v)} \hat\mu_t(e)\,\hat w_t(e)\Bigr)\, x^{\mathrm{new}}_t(k,v) .}
\end{equation*}
\new{Demand and congestion-aware storage aim for a small value of}
\begin{equation*}
\new{D + \beta K ,}
\end{equation*}
\new{where $\beta$ sets how much crowding counts against distance, and
demand storage has $\beta = 0$. They are greedy. They place one SKU at a
time on its best shelves with room, and they do not search for the smallest
value. Shelves that score the same go by the trip distance
$d^{\mathrm{trip}}_t(k,v)$, then by index.}

\subsection*{Appendix C, the arrivals (rewritten)}
\new{How much arrives is set by the orders alone, so it is the same under
every rule. Stock comes in refills of $Q$ units. A SKU gets one refill for
every $Q$ units ordered since its last refill, and units that do not make a
whole refill count towards the next one. So if 23 units of tea were ordered
since its last refill and $Q = 10$, two refills of tea arrive and 3 units
carry over. Formally, with $\nu_t(k)$ the units of SKU $k$ ordered since its
last refill, counted when their lines are released, the arrivals are}
\begin{equation*}
\new{y_t(k) = Q \left\lfloor \frac{\nu_t(k)}{Q} \right\rfloor ,}
\end{equation*}
\new{and $\nu_t(k) - y_t(k)$ carries over to the next update.}

\textit{The old text, which fills each SKU up to a level with a margin and a
floor, and its equation are removed. The notation table gains $\nu_t(k)$.}

\subsection*{Other places that name the crowding cost}
\new{The loop section, the scoring section and the Background refer to ``the
congestion term $K$'' of the storage objective. They would say ``the crowding
cost of congestion-aware storage (Section VI-A)'' instead, so the reader is
not sent to an equation in a later chapter.}

\clearpage
\section*{Part 4. The code}
For the thesis settings (\texttt{configs/step12}), checked in
\texttt{storage/arrival.py} and \texttt{environment/multi\_agent\_env.py}.
\begin{enumerate}[leftmargin=1.5em]
\item \textbf{$b_k$ is real.} The SKUs take 0.5, 1 or 2 units of room, in turn.
\item \textbf{Whole refills, unplaced refills, the leg weights, the ties and
the refill arrivals} already work as Parts~2 and~3 describe. With these
choices the thesis follows the code, so no rule changes and nothing has to
be run again.
\item \textbf{Only comments in the code are out of date.} The top comment of
\texttt{storage/arrival.py} still describes the trip from the nearest dock
and the old put-away score, and uses the old names (receiving vertex,
delivery vertex). These would be brought in line with the thesis, with no
change to what the code does. This goes in a small pull request on
\texttt{slap-mapd-coupling}, which you open.
\end{enumerate}

\subsection*{Questions for you}
\begin{enumerate}[leftmargin=1.5em]
\item Do Parts~2 and~3 go into the thesis as shown?
\item Should the code comments be updated as in Part~4, point~3?
\item Should the differences found be noted in the tracking issue \#103,
marked as fixed by the thesis side?
\end{enumerate}
\end{document}
